\section{Binding Flaws in Audited Applications}\label{sec:study}

We ask how common binding flaws are in practice, how severe they are, and
which patterns dominate (RQ1). Professional audits are the best available
record: each finding was found by an expert, argued in writing, and usually
fixed.

\subsection{Collecting audit findings}\label{sec:study-collect}

\para{Sources}
We enumerated 8,578 unique public audit reports from 68 audit sources: firms, contest
platforms and advisory databases, including the report archives of firms
that audit ZK systems (e.g., Veridise, Zellic, zkSecurity, Trail of Bits,
ABDK, OpenZeppelin, Consensys Diligence) and the public results of
Code4rena, Sherlock and Cantina.
% src: dataset/audit-sweep/merged/FUNNEL.md Overall (8,578 unique audits); all_reports.csv (68 auditor_or_platform values)
Some sources could not be fully enumerated during the collection window. We
list them with their coverage in the artifact.
\TODO{W6}{one-line coverage statement per missing source (Solodit, CertiK, yAcademy) from COVERAGE.md}

\para{Screening}
A report is in scope if it audits an application that verifies a ZK proof and
acts on it. 343 reports passed this screen. A stricter deduplication by
report URL gives 322 distinct reports, of which 311 publish their finding
text; 311 is the denominator for every per-report rate below.
% src: labels-v1/D2-TREND.md reconciliation (345 IN rows -> 322 URLs; 11 without public text -> 311)
 From them we logged 698 findings
that touch proof verification, public inputs, nullifiers, verifier keys or the
use of verified values (662 after merging findings shared across reports).
After screening, 300 candidates remained. We labeled all of them except
those triaged as likely not binding, plus a random sample of 14 of those to
check the triage: 264 items.
% src: merged/FUNNEL.md Overall (343 IN; 698 rows; 662 groups; 300 candidates); labeling/key/MANIFEST.md (264; inclusion reasons incl. triage-LN-sample 14)
As a completeness check, the sweep recovered 49 of 50 binding findings in an
independently assembled list of known cases.
% src: merged/RECALL-R2.md (49/50)

\para{Labeling}
Each candidate was labeled against a written guide that defines C1--C3, the
five patterns, the exclusions of Section~\ref{sec:definition}, and worked
examples. Each label records binding (Y/N/U), pattern, whether the issue is a
proof-system soundness bug, C1--C3 individually, the effect, the unbound
value, and confidence. Labels were proposed by an LLM and adjudicated by two
collaborating raters who discussed every item and fixed the final label, so
the labels are not independent annotations. The raters read the auditor's
text for all items, and for the 51 items with a public repository and
vulnerable commit (\emph{FULL} tier) they also recorded the intended policy
$\Pi$ of Section~\ref{sec:def-ledgers}. For the other items, $\Pi$ comes
from the auditor's description of the intended behavior. For every item used in the tool evaluation, a second
code-level pass checked the vulnerable location in the source.
\DEMO{An independent rater, blind to these labels, relabeled a stratified
sample of 30 items and agreed on the binding label with Cohen's
$\kappa=0.78$ and on the pattern with $\kappa=0.71$.}
\TODO{D1}{independent relabel of 30 stratified items; compute $\kappa$ for binding and pattern}
\DEMO{We then re-applied the policy rule to all 103 binding labels and to the
17 non-binding items labeled policy-bounded. Two labels changed: a capped
relayer fee became non-binding, and a reusable attestation that the
documentation describes as single-use became binding, which leaves the counts
below unchanged.}
\TODO{RV4}{re-check every P1/P3 label against the policy rule; list changed items}

Of the 264 labeled items, 4 duplicate an earlier item and 5 could not be
decided from the report (U), leaving 255 findings for analysis.
% src: dataset/labels-v1/SUMMARY.md (264 items; Y 103, N 156, U 5; dups 4)

\begin{table}[t]
\centering
\caption{From audit reports to binding findings.}
\label{tab:funnel}
\footnotesize
\begin{tabular}{@{}lr@{}}
\toprule
\textbf{Stage} & \textbf{Count} \\
\midrule
Audit reports enumerated (68 sources) & 8,578 \\
Reports on ZK-consuming applications & 343 \\
Proof-related findings logged (unique) & 698 (662) \\
Candidates after screening & 300 \\
Labeled items (FULL 51 / LIGHT 152 / PSS 61) & 264 \\
Decidable, non-duplicate findings & 255 \\
Binding flaws (with adjacent proof-system bugs) & 103 \\
\quad \textbf{core class (excl.\ proof-system soundness)} & \textbf{87} \\
\bottomrule
\end{tabular}
% src: merged/*.csv; labels-v1/SUMMARY.md; labels-v1/CHARACTERIZATION.md
\end{table}

\subsection{Prevalence}\label{sec:study-prev}

87 of the 255 labeled findings (34\%) are core binding flaws. With the 16
adjacent proof-system findings the total is 103 (40\%), which is 15\% of all
698 proof-related findings logged. The core findings occur in 43 of the 68
project families in the labeled set (44 with the adjacent ones), so they are
not concentrated in a few careless codebases.
% src: CHARACTERIZATION.md (103 of 255; 44 of 68 families); LABELS-v1.csv PSS!=Y: 87, 43 families
For contrast, 48 of the non-binding findings are circuit-internal soundness
issues (under-constrained signals with no effect on what the consumer is
told), which is the class that existing ZK tools target.
% src: CHARACTERIZATION.md (tagged circuit-soundness N: 48)

\finding{Binding flaws are a major class of ZK application vulnerabilities:
87 auditor-reported core instances, in 43 of the 68 project families we
studied.}

\subsection{Severity}\label{sec:study-sev}

Table~\ref{tab:severity} (Appendix~\ref{app:extra}) breaks the binding
findings down by pattern and auditor severity. 46 of the 87 core findings are rated Critical or High,
which is 78\% of the 59 core findings with a rating; with the adjacent
findings the share is 81\% (56 of 69). Only 7 are rated Low or
Informational. Findings without a rating are mostly contest and advisory
entries whose report gives no severity scale.
% src: CHARACTERIZATION.md pattern x severity (Critical 32, High 24, unrated 34, Low/Info 7); core from LABELS-v1.csv PSS!=Y: C 27, H 19, M 6, L 6, I 1, unrated 28
A binding flaw usually means an attacker obtains funds or rights that the
proof did not grant, so the severity is direct.



\finding{Binding flaws are severe: 78\% of rated core findings are Critical
or High.}

\subsection{Patterns}\label{sec:study-patterns}

Verifier integration (P4) is the largest pattern with 33 findings, but about
half of them (16) are proof-system soundness bugs that we report as adjacent.
Among the core class, unbound effects (P1, 23) and missing context (P2, 23)
dominate, followed by replay and aliasing (P3, 15). Interface mismatches (P5,
7) are rare, partly because a mismatch usually makes honest proofs fail, which
developers notice in testing.
Table~\ref{tab:examples} (Appendix~\ref{app:extra}) shows representative findings. Their variety is
the point. The unbound value can be a refund address, a bond receiver, a
prior state root, a chain ID, a zkVM program ID, or the position of an
attribute in the public-input array. Yet each is the same failure: the
consumer acts on something the statement does not fix.


\DEMO{Most P1 findings concern a recipient or fee that the contract takes from
calldata; most P2 findings concern a root or prior state that the consumer
does not check against its own storage.}
\TODO{D5}{tabulate the unbound value per pattern from LABELS-v1 \texttt{unbound\_value}}

\finding{Unbound recipients and amounts (P1) and missing context (P2) account
for over half of the core class. Both are the ZK analogues of signing too
little.}

\subsection{Trend and who finds them}\label{sec:study-trend}

Raw counts of core findings grow: 6 from 2021--2022, 11 from 2023, 19 from
2024, 32 from 2025, 6 from the partial year 2026, and 13 undated (87; with
adjacent findings 6, 13, 20, 39, 9 and 16, summing to 103). They also track
how many ZK audits were published.
% src: CHARACTERIZATION.md by report year; D2-TREND.md core table
The number of in-scope ZK audit reports grows too, from 60 in 2023 to 94 in
2025. Normalized by those reports, the share of reports with at least one
core finding is 9 of 60 in 2023 (0.15, 95\% Wilson CI 0.08--0.26), 9 of 58
in 2024 (0.16) and 21 of 94 in 2025 (0.22, 0.15--0.32). The change is not
significant (Cochran--Armitage over 2023--2025, $p=0.22$; $p=0.12$ with the
adjacent findings), so we do not claim a rising rate. Over all 311 reports
from 2021 to 2026, including undated ones, 55 (18\%) report at least one core
finding, and 63 (20\%) at least one binding finding of either kind.
% src: dataset/labels-v1/D2-TREND.md (analysis/d2_trend.py)
The 87 core findings come from 21 different auditors and advisory sources
(24 for all 103), and no source contributes more than 21. Binding is therefore a class that many
experienced auditors encounter, not the checklist item of one firm.
% src: CHARACTERIZATION.md by auditor (24 sources; max Veridise 21)

\para{Which layers are needed}
We tagged each core finding with the code a reviewer must read to decide it,
under one rule: the statement side counts as known when the audit report
states what the proof establishes. This rule favors single-layer tags,
because the auditor has already done the cross-layer reading and written it
down. \DEMO{The tags are LLM-proposed and author-checked on a stratified
sample of 15. Under this rule, 43 of the 87 core findings (49\%) can be decided from
the consumer contract alone, 20 (23\%) from the circuit alone, and 7 (8\%)
only by comparing the two. The remaining 17 (20\%) lie in verifier,
recursion or off-chain framework code. Under the alternative rule, a finding
is cross-layer whenever its fix adds the missing value to the statement; the
contract-only share then falls to 32\% and the cross-layer share rises to
25\%. The 7 is therefore a lower bound that depends on the report, not a
property of the code. Even the contract-only findings are not ordinary
contract analysis: deciding them still requires knowing \emph{which values
the statement fixes}, which contract analyzers do not model
(Section~\ref{sec:eval}).}
% src: dataset/labels-v1/D3-LAYERS.md, D3-layers.csv (AI-proposed tags)
\TODO{D3}{author spot-check of the 15-item stratified sample in D3-LAYERS.md; re-tag under amendment A3 if more than 2 of 15 disagree}

\para{How they were fixed}
\DEMO{Fixes for the core findings mostly mirror signature-misuse remedies:
56\% (49 of 87; 95\% CI 46--66\%) add the missing value to the statement,
bind a context, or make the replay key unique, and another 34\% pin the
verifier or program, check its result, or enforce a canonical encoding. The
ZK-specific residue is 6 findings fixed by constraining a value the circuit
already exposed. Most fix classes are inferred from the report; on the 21
findings whose fix is stated explicitly, 19 fall in these classes.}
% src: dataset/labels-v1/D4-FIXES.md, D4-fixes.csv (AI-proposed, report text only)
\TODO{D4}{author check of the AI-proposed fix classes (D4-FIXES.md)}

\finding{About one in six ZK application audits (55 of 311, 2021--2026)
reports a core binding flaw, found by many different auditors. Deciding one requires knowing what the
proof's statement fixes, whether the flaw sits in the contract, the circuit,
or between them. That reasoning is what Section~\ref{sec:design} automates.}
